\section{Results}\label{sec:results}

\subsection{Policy costs and the joint non-stress shortfall}\label{subsec:results_costs_shortfall}

Under the primary channel-based specification, no deterministic crossing occurs in the non-stress scenarios (low, mid, or frontier). The first reason is scale: the common denominator makes clear that the same policy label can imply very different commitments across settings. Table~\ref{tab:costs_results_clean} reports the six cost rows generated by the five instruments. Full UBI is the largest commitment everywhere, at 11.62--21.36 percent of endpoint GDP, followed by PBI at 4.85--7.97 percent. The relative position of the cheaper instruments is country-specific. Social pensions cost 0.56 percent of GDP in Peru but 3.10 percent in Chile, reflecting national benefit and eligibility rules. Ideal GMI costs only 0.15 percent in Chile, compared with 2.99 percent in Colombia; the loaded variant raises the cost in every country. A policy-cost gradient is therefore visible at broad scale, but there is no universal ordering among pensions, MUT, and GMI.

\begin{table}[H]
\centering
\small
\begin{singlespace}
\caption{Comparable gross policy costs (percent of endpoint GDP)}
\label{tab:costs_results_clean}
\begin{tabular}{@{}lrrrr@{}}
\toprule
Instrument & Chile & Colombia & Mexico & Peru \\
\midrule
Social pension (PEN) & 3.10 & 0.90 & 1.60 & 0.56 \\
Minimum universal transfer (MUT) & 2.90 & 4.75 & 3.95 & 5.34 \\
GMI, ideal targeting & 0.15 & 2.99 & 1.01 & 2.79 \\
GMI, loaded targeting & 0.22 & 4.48 & 1.52 & 4.18 \\
Partial basic income (PBI) & 4.85 & 7.43 & 5.91 & 7.97 \\
Full UBI benchmark & 11.62 & 18.98 & 15.81 & 21.36 \\
\bottomrule
\end{tabular}
\begin{minipage}{0.94\textwidth}\footnotesize
\emph{Note:} GMI costs are fixed aggregate endpoint costs; demographic instruments use the fixed endpoint eligible-share rule. Gross and net columns coincide because the primary comparison credits no existing-spending offset.
\end{minipage}
\end{singlespace}
\end{table}


The table is a financing comparison, not a ranking of social value. A low cost can reflect narrow eligibility, a low benefit, a small measured poverty gap, or ideal targeting; it does not establish adequacy or administrative desirability. Conversely, the UBI row is retained because its universal denominator makes it a useful fiscal stress benchmark, not because it is the preferred policy in every country. Keeping gross cost in the primary comparison also prevents assumed program replacement from mechanically improving $V$. Any credit for substitutable existing spending belongs to a labeled sensitivity rather than to the cross-instrument ranking.

The household extensions support using the aggregate GMI rows without making them appear exact. At the matched international poverty threshold, the primary CASEN estimate is 0.39 percent above Chile's aggregate GMI cost. The corresponding ENAHO estimate is 7.73 percent below Peru's aggregate cost. In both cases the loaded version inherits the same relative difference because it applies the same targeting-cost factor. These comparisons validate the order of magnitude while leaving the official four-country grid on its common aggregate basis.

The second reason is that the domestic gain and fiscal conversion are jointly too small. Across all 360 low, mid, and high cells, no value reaches the structural accounting threshold. The largest is $V^{gross}=0.177$, for Peru's social pension under the mid scenario and r4. The high frontier input does not overturn this result; as explained in Section~\ref{subsec:scenarios_regimes_clean}, its country bridge can place it below the mid input after conversion. The absence of a high-scenario crossing is therefore a result of the model-ready inputs, not an imposed monotonic scenario ordering.

The diagnostic inversions show why the shortfall should not be assigned to one margin. With the 10-percent prudential buffer, every mid-scenario accounting and debt-consistent cell requires $T^{req,H}>1$. The smallest requirement is 1.0941, for Chile's ideal GMI under r4, already beyond the domestic-translation frontier. Holding that cell's translation fixed instead, the minimum required AI-rent capture rate is 202.46 percent. This is 13.50 times the most permissive specified upper support of 15 percent under r4. Thus frontier translation alone does not close the gap, and neither does capture at the strongest specified reform setting. Panel B of Table~\ref{tab:frequency_threshold_hstar_clean} shows that the required domestic gain and gross marginal fiscal conversion also rise sharply for the corresponding Peruvian, Colombian, and Mexican instruments. The mid-scenario result is a joint shortfall in translated income and public capture. Because productive exposure uses a common LAC fallback, cross-country translation differences in the evaluated grid operate through preparedness and adoption; the design does not separately identify a country-specific productive-exposure gradient.

Some low-scenario values are negative because fixed administration, transition, and leakage deductions exceed the small revenue flow generated by the translated gain. This does not imply that AI reduces output in those cells; it means that the incremental fiscal-space numerator is negative after implementation costs. As the scenario strengthens, reform raises marginal conversion and the fixed deductions become less important relative to the gain. Yet the best mid cell remains far below coverage, so the result is not a knife-edge consequence of one small fixed-cost assumption.

\subsection{Stress crossings and the debt guardrail}\label{subsec:results_stress_debt}

Only the disruptive stress scenario produces deterministic accounting crossings, and the debt guardrail separates Chile's GMI result from Peru's pension result. Figure~\ref{fig:headline_feasibility_clean} displays the r0 grid on a common categorical scale. It makes three features visible at once: there is no crossing in the lower three scenario columns; both Chilean GMI cost variants cross under stress; and Peru's pension crossing is marginal relative to the prudential buffer.

\begin{figure}[H]
\centering
\includegraphics[width=\textwidth]{figures/headline_feasibility.pdf}
\caption{Primary feasibility grid under the status-quo fiscal regime. Each entry is deterministic $V^{gross}$ from the primary grid, rounded to two decimals. The black outline marks the structural accounting threshold $V\geq1$; the categorical fill distinguishes cells below zero, below one, between 1 and 1.10, and at or above the 10-percent prudential buffer.}
\label{fig:headline_feasibility_clean}
\end{figure}

Under r0, Chile's ideal GMI reaches $V=3.90$ and its loaded companion reaches $V=2.60$. Both pass the debt guardrail, and both remain debt-consistent across r0--r4. Combined reform raises the corresponding values to 5.80 and 3.87. The loaded result matters because it shows that the crossing is not created by ideal targeting alone. Its interpretation is nevertheless conditional: it requires persistence of the disruptive scenario that the design treats as a diagnostic stress test, not as a central macroeconomic projection.

Peru's social pension also crosses the structural accounting threshold under stress in every regime, with $V$ ranging from 1.02 to 1.42. Only r2 and r4 also clear $V=1.10$. More importantly, none closes the inherited primary-balance gap: every row is classified as accounting feasible but debt failed. The result is therefore fragile in a fiscal sense even where reform increases contemporaneous coverage. Colombia and Mexico never reach the accounting threshold. Their strongest stress cells are Colombia's pension under r4, at 0.91, and Mexico's ideal GMI under r4, at 0.80. Table~\ref{tab:deterministic_results_clean} collects these classifications.

The guardrail is doing distinct work here. It does not reprice the transfer or reduce the AI-induced revenue flow; it asks whether that flow can finance the policy while also covering the primary-balance correction inherited from the maintained debt path. A structural crossing can therefore remain informative as an accounting result and still be insufficient as a fiscal commitment. That distinction is why the Chile and Peru stress rows should not be grouped under a single feasibility label.

\begin{table}[H]
\centering
\small
\begin{singlespace}
\caption{Selected deterministic outcomes and the debt guardrail}
\label{tab:deterministic_results_clean}
\begin{tabular}{@{}>{\raggedright\arraybackslash}p{0.31\textwidth}>{\raggedright\arraybackslash}p{0.16\textwidth}r >{\raggedright\arraybackslash}p{0.34\textwidth}@{}}
\toprule
Scope or cell & Reported regime & $V^{gross}$ & Result class \\
\midrule
All low/mid/high cells & Best of 360 cells & 0.177 & No accounting crossing \\
Chile, GMI ideal, stress & r0 & 3.90 & Feasible, including debt guardrail \\
Chile, GMI loaded, stress & r0 & 2.60 & Feasible, including debt guardrail \\
Peru, social pension, stress & r0--r4 & 1.02--1.42 & Accounting crossing; debt guardrail fails \\
Colombia, best cell (PEN, stress) & r4 & 0.91 & No accounting crossing \\
Mexico, best cell (GMI ideal, stress) & r4 & 0.80 & No accounting crossing \\
\bottomrule
\end{tabular}
\begin{minipage}{0.94\textwidth}\footnotesize
\emph{Note:} Accounting crossings use the structural threshold $V\geq1$; Figure~\ref{fig:headline_feasibility_clean} separately identifies cells that also reach $V\geq1.10$. Stress is a diagnostic scenario, not a forecast.
\end{minipage}
\end{singlespace}
\end{table}


The deterministic evidence consequently supports a narrow hierarchy of findings. Chilean GMI is conditionally feasible only under diagnostic stress; Peru's pension has an accounting crossing without debt consistency; and neither Colombia nor Mexico has a crossing in the evaluated grid. Reform changes the distance to the threshold, but it does not turn a lower AI scenario into a feasible cell.

\subsection{Crossing frequencies, threshold requirements, and persistence}\label{subsec:results_mc_hstar}

Calibrated uncertainty reinforces the conditional nature of the deterministic crossings. In the primary channel Monte Carlo, the largest model-implied frequency over low, mid, and high scenarios is 2.36 percent, again for Chile's ideal GMI under mid and r4. Allowing the specified factor-dependence structures raises the maximum to 4.44 percent, without changing the absence of a deterministic non-stress crossing. These are frequencies across calibrated draws, not predictive probabilities; their scale is conditional on the parameter supports, distributional forms, and dependence design described in Section~\ref{subsec:uncertainty_clean}.

The historical reduced-form modes produce a deliberately different result. Historical increment ratios embed discretionary tax changes, commodity cycles, crisis recoveries, denominator noise, and shifts in revenue composition; they serve as reduced-form robustness benchmarks rather than as the preferred estimate of marginal capture from an AI-specific increment. For Chile's GMI under mid and r0, drawing gross capture from the historical distributions gives crossing frequencies of 66.06--91.76 percent for ideal GMI and 55.78--85.10 percent for loaded GMI. These modes are available only under r0 and replace the channel construction with a historical-capture draw; they do not show that the primary channel model predicts a likely mid-scenario crossing. Rather, the contrast identifies fiscal conversion as an empirically consequential modeling choice. The total-revenue variants occupy the upper part of these ranges, while the tax-based mixtures are lower.

Under stress, Chilean GMI has primary channel frequencies of 94.78--99.98 percent across its two cost variants and five regimes, with the debt-consistent frequency unchanged. Peru's pension is different: its accounting crossing frequency ranges from 41.02 to 80.50 percent, but the debt-consistent frequency is only 0.02--2.48 percent. The latter upper value occurs under r4; it is small, but not below one percent. Panel A of Table~\ref{tab:frequency_threshold_hstar_clean} keeps the primary, dependence, and historical objects separate.

This comparison also clarifies why a deterministic crossing is not a probability statement. The point calibration fixes one value for every primitive, whereas the simulation draws the scenario, translation, policy cost, and fiscal channels jointly within their stated supports. Chile's stress result remains well inside the crossing region across most draws. Peru's accounting result lies nearer the boundary, so its frequency is more responsive to the fiscal regime, and the debt correction moves almost the entire calibrated distribution back below the stronger requirement.

The inversion panel turns these crossing frequencies into requirements. Chile's ideal GMI under mid and r4 is the closest non-stress cell, while Peru's pension, Colombia's pension, and the Mexican GMI variants require progressively larger accumulated gains and gross fiscal conversion. In every listed row the required translation exceeds one, and the AI-rent inversion lies outside a feasible tax-rate domain. The exercise therefore reports how far each cell is from the boundary rather than treating every non-crossing as equivalent.

The persistence inversion gives the same comparison in years, but it is a post-baseline derived extension. At the 10-percent buffer, Chile's ideal GMI under stress requires six years of persistence under r0 and five under r4; the debt-consistent horizon is the same. Peru's pension has an accounting horizon of 8--11 years, but the debt-consistent requirement is 17--22 years. Colombia's pension reaches the accounting boundary in 13--19 years, while its debt-consistent counterpart is censored beyond 25 years. Mexico's ideal GMI requires 13--16 years and its loaded GMI 17--21 years on the accounting basis; both debt-consistent versions are likewise beyond 25 years. These are conditional persistence requirements, not calendar forecasts or expected adoption dates.

\begin{table}[H]
\centering
\footnotesize
\begin{singlespace}
\caption{Crossing frequencies, threshold requirements, and persistence horizons}
\label{tab:frequency_threshold_hstar_clean}
\textit{Panel A. Model-implied crossing frequencies across calibrated draws}\\[0.2em]
\begin{tabular}{@{}>{\raggedright\arraybackslash}p{0.48\textwidth}rr@{}}
\toprule
Mode and cell & Accounting (\%) & Debt-consistent (\%) \\
\midrule
Primary channel MC, maximum over low/mid/high & 2.36 & 2.36 \\
Factor-dependence modes, maximum over low/mid/high & 4.44 & 4.44 \\
Historical reduced form, Chile GMI ideal, mid r0 & 66.06--91.76 & 66.06--91.76 \\
Historical reduced form, Chile GMI loaded, mid r0 & 55.78--85.10 & 55.78--85.10 \\
Primary channel MC, Chile GMI, stress r0--r4 & 94.78--99.98 & 94.78--99.98 \\
Primary channel MC, Peru PEN, stress r0--r4 & 41.02--80.50 & 0.02--2.48 \\
\bottomrule
\end{tabular}
\\[0.7em]\textit{Panel B. Buffered mid-scenario inversions under r4}\\[0.2em]
\begin{tabular}{@{}lrrrr@{}}
\toprule
Country and instrument & $g^{AI,*}$ (\%) & $MFC^{gross,*}$ (\%) & $T^{req,H}$ & $\tau^{AI,req}$ (\%) \\
\midrule
Chile, GMI ideal & 4.75 & 16.72 & 1.09 & 202.46 \\
Peru, social pension & 7.53 & 30.31 & 1.71 & 671.42 \\
Colombia, social pension & 10.62 & 45.23 & 2.39 & 1149.88 \\
Mexico, GMI ideal & 12.55 & 52.38 & 2.80 & 1364.07 \\
Mexico, GMI loaded & 17.31 & 72.23 & 3.79 & 2025.71 \\
\bottomrule
\end{tabular}
\\[0.7em]\textit{Panel C. Required persistence $H^*$ under stress, $\xi=0.10$}\\[0.2em]
\begin{tabular}{@{}>{\raggedright\arraybackslash}p{0.47\textwidth}rr@{}}
\toprule
Country and instrument & Accounting $H^*$ & Debt-consistent $H^*$ \\
\midrule
Chile, GMI ideal (r0 / r4) & 6 / 5 & 6 / 5 \\
Peru, social pension (r0--r4) & 8--11 & 17--22 \\
Colombia, social pension (r0--r4) & 13--19 & >25 \\
Mexico, GMI ideal (r0--r4) & 13--16 & >25 \\
Mexico, GMI loaded (r0--r4) & 17--21 & >25 \\
\bottomrule
\end{tabular}
\begin{minipage}{0.94\textwidth}\footnotesize
\emph{Note:} Frequencies use \texttt{prob\_basis=all\_draw}. Historical reduced-form modes are alternative capture specifications available only under r0. Panel B reports accounting inversions at $\xi=0.10$; $g^{AI,*}$ is an accumulated gain. $H^*$ is a post-baseline derived extension and is censored after 25 years.
\end{minipage}
\end{singlespace}
\end{table}


\subsection{Diagnostic and robustness}\label{subsec:results_diagnostics}

The ICT placebo does not reproduce the headline pattern, while the exposure-rescaling diagnostic shows that the strongest Chilean stress crossing is not specific to the productive-exposure mechanism. The placebo produces no crossing in 120 country--policy--regime cells; its maximum is $V^{placebo}=0.750$. Ordinary ICT-era digitalization therefore does not reproduce the threshold pattern when paired with era-consistent fiscal capture and policy-cost proxies.

The reduced-form diagnostic multiplies the primary ratio by $E^{auto}/E^{prod}=0.318$. It removes nine of the 15 structural crossings, including every Peruvian pension crossing. Six survive: Chile's ideal GMI under all regimes and its loaded GMI under r4. Because the ideal-GMI stress crossing remains after productive exposure is replaced by automation-risk exposure, the exposure-rescaling diagnostic does not clear the crossing criterion. It narrows the mechanism claim without changing the accounting comparison.

Because the diagnostic rescales $V$ rather than re-estimating the translation block, it is a reduced-form sensitivity rather than a complete negative-control design. A lost crossing shows dependence on productive rather than automation exposure in the maintained reduced form. A surviving crossing shows only that the original stress gain and low policy cost remain sufficient under the alternative scaling; it does not establish an automation-based fiscal mechanism.

Channel ablations further distinguish the two crossing types. Chile's ideal-GMI stress cell remains above one when labor, capital, consumption, or AI-rent capture is removed individually, so no single fiscal channel creates that crossing. Peru's pension under r4 falls below one when labor or consumption is removed, but survives without capital or AI rents. These diagnostic values use the behavioral-erosion companion, which is why their recomputed baseline is not substituted for the deterministic value in Table~\ref{tab:deterministic_results_clean}. The adoption-side informality ablation is inert in the primary baseline because the adoption ceiling is not binding; fiscal informality remains embedded in observed effective rates and cannot be separately removed without a new counterfactual.

The available leave-one-source-out exercise is narrower. Three like-for-like OECD-to-GRD revenue substitutions preserve both class and rank. Comparable exclusions for WDI, PIP, and AIPI are unavailable in the archived inputs and are reported as such rather than filled by an imputed alternative. The result supports the evaluated revenue substitution, not a claim of source invariance across the entire model.

\begin{table}[H]
\centering
\footnotesize
\begin{singlespace}
\caption{Diagnostic and robustness results}
\label{tab:diagnostics_clean}
\begin{tabular}{@{}>{\raggedright\arraybackslash}p{0.25\textwidth}>{\raggedright\arraybackslash}p{0.67\textwidth}@{}}
\toprule
Exercise & Result \\
\midrule
ICT temporal placebo & No crossing in 120 cells; maximum $V^{placebo}=0.750$. \\
Reduced-form exposure-rescaling diagnostic & 6 of 15 structural crossings survive after multiplication by 0.318; the diagnostic does not clear the crossing criterion. \\
Channel ablation & Chile GMI ideal remains above one after any single channel is removed; Peru PEN falls below one without labor or consumption but not without capital or AI rents. \\
Leave-one-source-out & The 3 available OECD--GRD substitutions preserve class and rank; 3 core-source exclusions are unavailable in the archived inputs. \\
\bottomrule
\end{tabular}
\begin{minipage}{0.94\textwidth}\footnotesize
\emph{Note:} Channel-ablation values use the diagnostic recomputation with the behavioral-erosion companion.
\end{minipage}
\end{singlespace}
\end{table}


Section~\ref{sec:discussion} now interprets the accounting and diagnostic results for policy design.
